\paragraph{结论。}\label{sub:conclusions}
我们不再试图从\trPicard--\trLefschetz{}数据几何地恢复诸 $X^k$，而是在\trFukayaCategory{}的框架中代数地执行一个类似的构造。结果表明，上面遇到的那些困难在这一层面上消失了（对照\cref{thm:matching-isotopy-role-remark}）。为使行文简单，我们把 $\Psi^k$ 与 $\varpi^k$ 当作真的是\trExact\trLefschetzFibration{}与\trBifibration{}来对待。这是不成问题的：$\Psi^k$ 的情形已在\cref{sub:lefschetz-pencils}中解释过（对 $\varpi^k$ 的讨论将完全类似）。于是，我们有两列\trFukayaCategory{}$\scr{F}(X^k)$ 与 $\scr{F}(\Psi^k)$。对每个 $0<k\le n$，设 $V^k_j\subset X^k$ 是与 $\gamma^{k+1}_j$ 相伴的\trVanishingCycle，按构造它也是与 $\mu^k_j$ 相伴的\trMatchingCycle（见\cref{thm:vanishing-as-matching-cycle}）。我们照常把它们做成\trLagrange\trBrane，并把所得的集合记作 $\boldsymbol{V}^{k,\#}$。写
\begin{equation*}
\begin{aligned}
\scr{A}^k&=\scr{F}(\boldsymbol{V}^{k,\#}),\\
\scr{A}^{k,\rightarrow}&=\scr{F}^{\rightarrow}(\boldsymbol{V}^{k,\#});
\end{aligned}
\end{equation*}
具体地说，$\scr{A}^k$ 是由我们这个\trVanishingCycle{}集合构成的 $\scr{F}(X^k)$ 的全 \trAinfty-子\trCategory，而 $\scr{A}^{k,\rightarrow}$ 是与之相伴的\trDirected{}\trAinfty-子\trCategory。

我们算法的开端与前面相同，即从\trTopology{}上重构 $X^n$ 连同 $\mathcal{K}_{X^n}$ 的\trTrivialization，以及\trVanishingCycle{}$\boldsymbol{V}^n$。在此基础上，$\scr{A}^{n,\rightarrow}$ 可以按 §13 中所解释的那样组合地计算出来。迭代步骤如下。假设给了一个\trQuasiisomorphic{}于 $\scr{A}^{k+1,\rightarrow}$ 的\trDirected{}\trAinfty-\trCategory{}$\scr{B}^{k+1,\rightarrow}$。依\cref{thm:full-embedding-functor}与\cref{thm:tw-directed-quasiequiv}，我们得到一个\trCohomology{}\trFullyFaithful{}的 \trAFunctor
\begin{equation}\label{eq:fukaya-functor-composition}
\scr{F}(X^k)\longrightarrow\scr{F}(\Psi^k)\longrightarrow Tw\scr{A}^{k+1,\rightarrow}.
\end{equation}
如\cref{thm:matching-cycles-computable}的证明中那样，我们可以确定任何\trMatchingCycle{}在 \eqref{eq:fukaya-functor-composition} 之下的像，以及 $Tw\scr{B}^{k+1,\rightarrow}$ 中相应的\trObject——或许相差\trGrading{}的一次\trSuspension（那份含糊归根结底无关紧要；为简单起见，让我们假定诸\trGrading{}已经选对）。于是，$\scr{A}^k$ \trQuasiisomorphic{}于一个被显式给出的全 \trAinfty-子\trCategory{}$\scr{B}^k\subset Tw\scr{B}^{k+1,\rightarrow}$。如果 $k>1$，我们\emph{丢弃一部分信息}，只保留 $\scr{B}^{k,\rightarrow}$——它\trQuasiisomorphic{}于 $\scr{A}^{k,\rightarrow}$——然后带着下一个更高的维数继续。以这种方式到达 $k=1$ 时，我们手中是一个\trQuasiisomorphic{}于 $\scr{A}^1$ 的 \trAinfty-\trCategory{}$\scr{B}^1$。最后，我们还可以再走半步，把 \eqref{eq:fukaya-functor-composition} 再应用一次，这表明 $\scr{F}(X)$ 本身\trQuasiequivalent{}于 $Tw\scr{B}^{1,\rightarrow}$ 的某个（\trApriori{}未知的）全 \trAinfty-子\trCategory。

\begin{proof}[定理 B 在 $n\ge 1$ 时的证明]
鉴于上面的论证，这现在已经显而易见，此时所考虑的 \trAinfty-\trCategory{}是 $Tw\scr{B}^{1,\rightarrow}$。
\end{proof}
\begin{proof}[推论 C 在 $n\ge 1$ 时的证明]
依\cref{thm:directed-quasiisomorphic-dg}，我们可以把 $\scr{B}^{1,\rightarrow}$ 换成一个与之\trQuasiisomorphic{}的\trDirected{}\trAinfty-\trCategory{}$\scr{C}^{\rightarrow}$，后者实际上是一个 dg \trCategory。取\trCohomology{}\trFullyFaithful{}的 \trAFunctor{}$\scr{F}(X)\to Tw(\scr{C}^{\rightarrow})$。在必要时为既有的\trObject{}引入若干复本之后，可假设这一\trFunctor{}在\trObject{}上是\trInjective{}的，于是 $\scr{F}(X)$ 成为\trQuasiisomorphic{}于它的\trImage{}——那是一个 $\hom$ 空间为有限维的 dg \trCategory。
\end{proof}
\begin{proof}[定理 A 的证明]
依\cref{thm:vanishing-cycles-split-generate}，我们知道诸\trVanishingCycle{}\trSplitGeneration{}$\Pi\,Tw\,\scr{F}(X^1)$。依\cref{thm:split-generated-equivalence}与\cref{thm:split-closure-existence}，$\Pi\,Tw\,\scr{F}(X^1)$ \trQuasiequivalent{}于 $\Pi\,Tw\,\scr{A}^1$，从而也\trQuasiequivalent{}于 $\Pi\,Tw\,\scr{B}^1$。
\end{proof}

细心的读者会注意到，该算法在最低的维数并不给出全部信息（我们可以轻易确定 $\scr{F}(X^n)$ 的任何\trDirected{}\trAinfty-子\trCategory，但这\trApriori{}对 $\scr{F}(X^n)$ 本身行不通）。部分原因在于，我们没有从真正平凡的零维情形开始这一归纳过程；而这可以按如下方式补救。取 $\Psi^n\colon X^n\to\CC$，并按 §18 的程序定义 $\scr{F}(\Psi^n)$：即取沿一根\trSmooth\trFiber{}分歧的\trDoubleCover{}$\widetilde{X}^n\to X^n$，并考虑 $\ZZ/2$-不变的\trLagrange\trSubmanifold。特别地，我们在 $\widetilde{X}^n$ 中有 (B) 型\trEquivariant\trBrane，它们是与诸\trVanishingPath{}$\gamma^n_j$ 相伴的那些\trLefschetzThimble{}的\trPreimage。细察\cref{thm:fibration-exceptional-collection}的证明即可看出，它们构成 $\scr{F}(\Psi^n)$ 中的一个\trExceptionalCollection。于是，相伴的全 \trAinfty-子\trCategory{}与\trDirected{}\trAinfty-子\trCategory{}是\trQuasiisomorphic{}的。另一方面，该\trDirected{}\trAinfty-子\trCategory{}适用于 §13 的组合方法，因为在定义相关的\trFloerCochainGroup{}与 $\mu^d$ 操作时，我们可以使用标准\trComplexStructure{}与零\trInhomogeneous{}项（\cref{thm:floer-datum-assumption} 仍将得到满足）。在这一特殊局面下，一切\trPseudoholomorphic\trPolygon{}实际上都将是常值的；其\trRegularity{}由\cref{thm:winding-injectivity}得出，一如\cref{thm:regularity-fixed-set}与\cref{thm:regularity-polygons}的证明中那样。把所得的\trDirected{}\trAinfty-\trCategory{}记作 $\scr{B}^{n+1,\rightarrow}$。它随之带来一个\trCohomology{}\trFullyFaithful{}的\trFunctor
\begin{equation}\label{eq:fukaya-functor-zero-dimension}
\scr{F}(X^n)\longrightarrow Tw\scr{B}^{n+1,\rightarrow}.
\end{equation}
利用它，可以把定理 B 与推论 C 的证明延展到 $n=0$（此时 $X^n=X$）的情形。

\begin{remark}\label{thm:zero-dimensional-extension}
为了类似地把定理 A 也延展过去，需要确定 $\Psi^{n-1}$ 的诸\trVanishingCycle{}在 \eqref{eq:fukaya-functor-zero-dimension} 之下的像。不巧，这里有一个细小而棘手之处，它在更高的维数没有出现过：在\cref{thm:matching-l-cone}的类似物中，现在有一个二维的零次\trMorphism{}空间，而锥必须对该空间中正确的那个元素来取（它对应于\trVanishingCycle{}上的非平凡 \trPinStructure）。
\end{remark}
